\begin{textsample}{POS Dim 8 – human – Score 2.00 – rj/rj\_ef\_matematica\_2\_p1.txt}  \label{ex:f8_pos_009}
No Ensino Fundamental, a área de Matemática está presente articulando seus \textbf{eixos} temáticos.
Desta forma a BNCC propõe o ensino da Matemática dividido em Unidades Temáticas, com características de semelhança para um ensinamento de progressão contínua.
Para o desenvolvimento das habilidades \textbf{prevista} na Base Nacional Comum Curricular, é imprescindível levar em consideração os conhecimentos matemáticos e as experiências do estudante.
A aprendizagem está relacionada à compreensão de sentido dos objetos matemáticos.
Esses sentidos são consequências das conexões que os estudantes formam entre os objetos e seu cotidiano, entre eles e os diferentes \textbf{temas} matemáticos e, por fim, entre eles e os demais componentes curriculares.
A Matemática é uma peça importante na construção da cidadania, nos conhecimentos científicos e recursos tecnológicos.
O seu ensino deve ser de construção do pensamento lógico-matemático, despertando o espírito investigativo, procurando desenvolver nos alunos competências para compreender e transformar a \textbf{realidade}, de maneira prazerosa.

% matched lemmas: eixo, prever, realidade, tema
\end{textsample}
